
\documentclass{article}
\usepackage{colortbl}
\usepackage{makecell}
\usepackage{multirow}
\usepackage{supertabular}

\begin{document}

\newcounter{utterance}

\centering \large Interaction Transcript for game `referencegame', experiment `line\_grids\_rows', episode 3 with reflexion{-}Qwen3.5{-}4B.
\vspace{24pt}

{ \footnotesize  \setcounter{utterance}{1}
\setlength{\tabcolsep}{0pt}
\begin{supertabular}{c@{$\;$}|p{.15\linewidth}@{}p{.15\linewidth}p{.15\linewidth}p{.15\linewidth}p{.15\linewidth}p{.15\linewidth}}
    \# & $\;$A & \multicolumn{4}{c}{Game Master} & $\;\:$B\\
    \hline

    \theutterance \stepcounter{utterance}  
    & & \multicolumn{4}{p{0.6\linewidth}}{
        \cellcolor[rgb]{0.9,0.9,0.9}{
            \makecell[{{p{\linewidth}}}]{
                \texttt{\tiny{[P1$\langle$GM]}}
                \texttt{You have attempted this same task before and did not succeed. These are the plans you wrote after those attempts.} \\
\texttt{Trial \#1: New plan: My previous expression was too generic and failed to distinguish the target grid from the distractors because it only described the target's appearance without explicitly referencing the specific difference from the others. The target has Xs in columns 3, 4, and 5, while Distractor 1 has Xs in columns 4 and 5, and Distractor 2 has Xs in columns 1 and 4. Next time, I will construct the expression by explicitly stating the unique position of the filled cells relative to the empty ones, specifically mentioning that the leftmost two columns are empty while the rightmost three are filled, ensuring this description uniquely identifies the target among the three options.} \\
\\ 
\texttt{Trial \#2: New plan: My previous expression failed because it described the target grid's general layout but did not explicitly contrast it with the specific patterns of the distractors. Distractor 1 has Xs only in the rightmost two columns, while Distractor 2 has Xs in the leftmost column and the rightmost two columns. To succeed, I must construct an expression that explicitly states the target grid is the one where the **leftmost two columns are completely empty** and the **rightmost three columns are filled with Xs**, thereby distinguishing it from the other two grids which have Xs in the leftmost column.} \\
\\ 
\\ 
\texttt{You are given three grids, where each of them is 5 by 5 in size.} \\
\texttt{Grids have empty cells marked with "▢" and filled cells marked with "X".} \\
\texttt{Your task is to generate a referring expression that best describes the target grid while distinguishing it from the two other distractor grids.} \\
\texttt{The first grid is the target grid, and the following two grids are the distractors.} \\
\\ 
\texttt{Target grid:} \\
\\ 
\texttt{□ □ X X X} \\
\texttt{□ □ X X X} \\
\texttt{□ □ X X X} \\
\texttt{□ □ X X X} \\
\texttt{□ □ X X X} \\
\\ 
\texttt{Distractor grid 1:} \\
\\ 
\texttt{□ □ □ X X} \\
\texttt{□ □ □ X X} \\
\texttt{□ □ □ X X} \\
\texttt{□ □ □ X X} \\
\texttt{□ □ □ X X} \\
\\ 
\texttt{Distractor grid 2:} \\
\\ 
\texttt{X □ □ X X} \\
\texttt{X □ □ X X} \\
\texttt{X □ □ X X} \\
\texttt{X □ □ X X} \\
\texttt{X □ □ X X} \\
\\ 
\texttt{Instruction: Describe the target grid.} \\
\texttt{Generate the referring expression starting with the tag "Expression: " for the given target grid. Omit any other text.} \\
            }
        }
    }
    & \\ \\

    \theutterance \stepcounter{utterance}  
    & \multicolumn{4}{p{0.6\linewidth}}{
        \cellcolor[rgb]{0.9,0.9,0.9}{
            \makecell[{{p{\linewidth}}}]{
                \texttt{\tiny{[P1$\rangle$GM]}}
                \texttt{Expression: The grid where the leftmost two columns are completely empty and the rightmost three columns are filled with Xs.} \\
            }
        }
    }
    & & \\ \\

    \theutterance \stepcounter{utterance}  
    & & & \multicolumn{2}{p{0.3\linewidth}}{
        \cellcolor[rgb]{0.9,0.9,0.9}{
            \makecell[{{p{\linewidth}}}]{
                \texttt{\tiny{[GM$|$GM]}}
                \texttt{Expression: The grid where the leftmost two columns are completely empty and the rightmost three columns are filled with Xs.} \\
            }
        }
    }
    & & \\ \\

    \theutterance \stepcounter{utterance}  
    & & \multicolumn{4}{p{0.6\linewidth}}{
        \cellcolor[rgb]{0.9,0.9,0.9}{
            \makecell[{{p{\linewidth}}}]{
                \texttt{\tiny{[GM$\rangle$P2]}}
                \texttt{You have attempted this same task before and did not succeed. These are the plans you wrote after those attempts.} \\
\texttt{Trial \#1: New plan: I misidentified the grid by focusing on the rightmost columns without verifying the leftmost columns. The expression requires the leftmost two columns to be empty (▢), which only the Second grid satisfies, but I likely confused the visual pattern or miscounted the columns. Next time, I will explicitly check the leftmost two columns first to ensure they are empty, then verify the rightmost three are filled with X, and select the Second grid.} \\
\\ 
\texttt{Trial \#2: New plan: I misidentified the grid because I focused on the "rightmost three columns" condition but failed to correctly verify the "leftmost two columns" condition. Upon closer inspection, the Second grid has Xs in the rightmost three columns but also has Xs in the leftmost two columns (specifically, the first two columns are empty in the description, but the grid shows Xs in the first two columns? Wait, let me re{-}read the grids.} \\
\\ 
\texttt{Grid 1: Left 2 empty, Right 3 X.} \\
\texttt{Grid 2: Left 3 empty, Right 2 X.} \\
\texttt{Grid 3: Left 2 X, Right 3 empty.} \\
\\ 
\texttt{The expression says: "Xs filling the rightmost three columns and empty cells in the leftmost two columns."} \\
\texttt{This matches Grid 1 exactly.} \\
\texttt{Grid 2 has Xs in the rightmost *two* columns.} \\
\texttt{Grid 3 has Xs in the leftmost *two* columns.} \\
\\ 
\texttt{My mistake was selecting "Second" instead of "First". I likely miscounted the number of Xs in the rightmost columns of the Second grid or confused the column indices. Next time, I will explicitly count the number of Xs in the rightmost three columns and the number of empty cells in the leftmost two columns for each grid before selecting. I will verify that Grid 1 has exactly 3 Xs on the right and 2 empty on the left, while Grid 2 has only 2 Xs on the right.} \\
\\ 
\texttt{New plan: I selected the wrong grid because I miscounted the number of filled columns in the Second grid versus the First grid. The expression requires the rightmost *three* columns to be filled with Xs. Grid 1 satisfies this (3 Xs on the right), while Grid 2 only has 2 Xs on the right. I will carefully count the Xs in the rightmost three columns and the empty cells in the leftmost two columns for each grid to ensure an exact match before answering.} \\
\\ 
\\ 
\texttt{You are given three grids, where each of them is 5 by 5 in size.} \\
\texttt{Grids have empty cells marked with "▢" and filled cells marked with "X".} \\
\texttt{You are also given a referring expression that describes one of the given grids.} \\
\texttt{Your task is to select the grid that matches the given referring expression.} \\
\texttt{Generate only the number (in text) of the grid that matches the given expression by selecting first, second, or third.} \\
\\ 
\texttt{First:} \\
\\ 
\texttt{□ □ X X X} \\
\texttt{□ □ X X X} \\
\texttt{□ □ X X X} \\
\texttt{□ □ X X X} \\
\texttt{□ □ X X X} \\
\\ 
\texttt{Second:} \\
\\ 
\texttt{□ □ □ X X} \\
\texttt{□ □ □ X X} \\
\texttt{□ □ □ X X} \\
\texttt{□ □ □ X X} \\
\texttt{□ □ □ X X} \\
\\ 
\texttt{Third:} \\
\\ 
\texttt{X □ □ X X} \\
\texttt{X □ □ X X} \\
\texttt{X □ □ X X} \\
\texttt{X □ □ X X} \\
\texttt{X □ □ X X} \\
\\ 
\texttt{Expression: The grid where the leftmost two columns are completely empty and the rightmost three columns are filled with Xs.} \\
\texttt{Question: Which grid does the expression refer to?} \\
\texttt{Start with the tag "Answer: ", followed by your selection. Omit any other text.} \\
            }
        }
    }
    & \\ \\

    \theutterance \stepcounter{utterance}  
    & & & \multicolumn{4}{p{0.6\linewidth}}{
        \cellcolor[rgb]{0.9,0.9,0.9}{
            \makecell[{{p{\linewidth}}}]{
                \texttt{\tiny{[GM$\langle$P2]}}
                \texttt{Answer: First} \\
            }
        }
    }
     \\ \\

    \theutterance \stepcounter{utterance}  
    & & & \multicolumn{2}{p{0.3\linewidth}}{
        \cellcolor[rgb]{0.9,0.9,0.9}{
            \makecell[{{p{\linewidth}}}]{
                \texttt{\tiny{[GM$|$GM]}}
                \texttt{first} \\
            }
        }
    }
    & & \\ \\

\end{supertabular}
}

\end{document}
